\chapter*{Önsöz}
\markboth{Önsöz}{Önsöz}
\addcontentsline{toc}{chapter}{Önsöz}

Bu kitap, IMO301 İstatistik ile PDR209 Temel İstatistik dersleri için
hazırlanmış iki ayrı çalışmanın ortak bir bilimsel ve pedagojik omurgada
birleştirilmesiyle oluşturulmuştur. Amaç, aynı kavramları iki kez anlatmak
değil; matematik eğitiminin kuramsal açıklama gücü ile PDR araştırmalarının
bağlamsal veri okuryazarlığını tek bir ilerleyişte buluşturmaktır.

Kitabın ortak çekirdeği araştırma sorusu, veri üretimi, betimleme, veri
kalitesi, eksik veri, örnekleme, belirsizlik, tahmin ve istatistiksel çıkarım
zinciridir. Bu çekirdek üzerinde
hipotez testleri, etki büyüklüğü, grup karşılaştırmaları, kategorik veri,
korelasyon, regresyon, parametrik olmayan yöntemler ve ölçek puanlarının
güvenirliği geliştirilir. Son bölümlerde çoklu regresyon ile bütünleştirici veri analizi laboratuvarı,
okuyucuyu tek bir yöntem uygulamaktan araştırmanın tamamını yönetmeye taşır.

Ortak konularda iki eski metinden daha kapsamlı ve uygulama bakımından daha
güçlü olan anlatım ana metin olarak seçilmiştir. Diğer metindeki farklı
örnekler, matematiksel ayrıntılar, uyarılar ve uygulamalar uygun bölümlere
aktarılmış veya kaynak havuzunda korunmuştur. Böylece yinelenen açıklamalar
çoğaltılmadan iki projenin özgün katkıları kaybolmamıştır.

Her bölümde şu soruların birlikte düşünülmesi önerilir: Araştırma sorusu
nedir? Veri nasıl üretilmiştir? Hedef parametre nedir? Yöntemin koşulları
sağlanıyor mu? Belirsizlik ne kadardır? Sonuç hangi evrene ve hangi bağlama
genellenebilir? Yazılım çıktısı ancak bu sorular yanıtlandıktan sonra anlam
kazanır.

Kitap doğrusal biçimde okunabileceği gibi, izleyen izlence rotası yardımıyla
IMO301 veya PDR209 dersinin haftalık kapsamına göre de kullanılabilir.

\begin{flushright}
İbrahim Güney\\
İstanbul, 2026
\end{flushright}